\section{从 Web 信息检索到检索增强生成的研究综述}

\subsection{关键词检索与 Web 索引}

早期 Web 搜索首先要解决的问题，是如何持续获取分散在网络中的网页，并让用户能够在可接受的时间内从大规模文档集合中找到相关内容。搜索引擎通常从一组种子地址出发，由爬虫下载页面、解析超链接并发现新的地址，同时进行去重、更新检测和异常处理。与结构相对稳定的封闭文档库相比，Web 具有规模大、变化快、内容异构且发布质量不可控等特点，因而网页抓取、索引更新和查询处理必须被设计为可扩展的系统。早期 Google 搜索引擎的公开设计已经将高效爬取、全文与超链接索引以及大规模查询处理视为相互关联的工程问题\parencite{brin1998anatomy}。

网页进入文档集合后，需要经过正文提取、分词、大小写或词形规范化等预处理，再被转换为可检索的表示。倒排索引是其中最关键的数据结构：系统为每个词项维护一个倒排表，记录包含该词项的文档编号，并可附加词频、出现位置等信息。查询到来时，系统只需访问查询词对应的倒排表，而无须逐篇扫描全部网页；倒排表还可通过跳表、压缩和分布式切分进一步提高处理效率\parencite{manning2008ir}。这一结构奠定了现代搜索系统的效率基础，即使在稠密检索出现之后，稀疏倒排索引仍因速度快、成本低和词项匹配可解释而被广泛使用。

在最基本的布尔检索中，文档被表示为词项是否出现的二值集合，系统利用 AND、OR 和 NOT 等逻辑运算合并倒排表。该方法规则清晰，却只能判断文档是否满足条件，难以给大量匹配结果确定优先顺序。排序检索因此引入词频与逆文档频率：词项在某篇文档中出现得越多，通常越能描述该文档；在整个集合中出现得越少，其区分能力通常越强。常见的对数化 TF--IDF 权重可写为
\begin{equation}
  w_{t,d}=\bigl(1+\log \mathrm{tf}_{t,d}\bigr)
  \log\frac{N}{\mathrm{df}_{t}},
  \label{eq:tfidf}
\end{equation}
其中，$\mathrm{tf}_{t,d}$ 表示词项 $t$ 在文档 $d$ 中的出现次数，$\mathrm{df}_{t}$ 表示包含该词项的文档数，$N$ 为文档总数。查询和文档由此可以表示为带权向量，并通过余弦相似度等方法计算相关性\parencite{manning2008ir}。

关键词检索及其索引机制实现了从“遍历网页”到“按词项定位文档”的转变，但其相关性判断主要依赖可观察的词面重合。当查询与文档使用同义表达时，系统可能漏掉语义相关结果；当同一词语具有多种含义时，又可能召回大量无关页面。TF--IDF 的词袋表示还弱化了词序、上下文和用户意图，面对自然语言问题、长尾查询以及关键词堆砌等情况时能力有限。因此，倒排索引解决了大规模检索的效率问题，却没有完全解决“哪些结果更有用”的问题，研究重点随之转向更精细的候选召回与相关性排序。

\subsection{候选召回与相关性排序}

倒排索引能够迅速定位含有查询词的网页，却可能返回数量庞大的匹配集合，而复杂的相关性计算又难以直接施加于全部网页。搜索系统因此逐渐形成“候选召回—相关性排序”的分阶段架构：召回阶段利用低成本方法从全量索引中筛选规模较小的候选集，首先保证潜在相关文档不被遗漏；排序阶段再使用更丰富的信号估计每个候选结果的价值。两阶段设计实质上是在效率与效果之间进行分工，前者更关注覆盖率，后者更关注最靠前结果的质量。若相关文档在召回阶段已经被过滤，后续排序模型再准确也无法将其恢复；反之，盲目扩大候选集又会增加延迟和计算成本。

BM25 是传统稀疏召回与排序中具有代表性的概率相关性模型。对于查询 $q$ 和文档 $d$，其常见形式为
\begin{equation}
  \operatorname{score}(q,d)
  =\sum_{t\in q}\operatorname{IDF}(t)
  \frac{\mathrm{tf}_{t,d}(k_1+1)}
  {\mathrm{tf}_{t,d}+k_1\left(1-b+b\frac{|d|}{\operatorname{avgdl}}\right)},
  \label{eq:bm25}
\end{equation}
其中，$|d|$ 与 $\operatorname{avgdl}$ 分别表示当前文档长度和集合平均文档长度，$k_1$ 控制词频增益的饱和速度，$b$ 控制长度归一化程度。与简单 TF--IDF 相比，BM25 不再假定词项重复出现会带来线性增益，并通过长度归一化减轻长文档因包含更多词语而天然占优的问题\parencite{robertson2009bm25}。它计算开销较低、无需大规模标注数据，至今仍适合作为关键词召回方法和更复杂模型的基线，但无法从根本上解决查询与文档之间没有共同词项时的语义匹配问题。

Web 搜索还能够利用普通文本集合所不具备的超链接结构。链接分析将网页及其链接关系表示为图，把其他页面的引用视为对目标页面价值的一种信号。PageRank 的基本思想是：来自重要页面的链接应当具有更大的权重，页面重要性因而通过链接图上的递归计算获得\parencite{brin1998anatomy}。这类与具体查询相对独立的权威性信号，可以补充关键词相关性并抑制部分低质量页面。然而，链接数量不等同于内容与当前需求相关，热门或存在时间较长的网站更容易积累链接，人工链接网络也可能操纵得分。因此，链接分析适合作为排序依据之一，而不能代替对查询与内容关系的判断。

随着可用特征和交互数据增加，学习排序（Learning to Rank, LTR）开始用监督学习代替固定的人工加权规则。模型可以综合 BM25 得分、标题与正文匹配程度、链接权威性、内容新鲜度以及用户点击等特征，根据人工相关性标注或隐式反馈学习排序函数。按照学习目标的不同，相关方法通常分为逐点、成对和列表三类：逐点方法独立预测每篇文档的相关等级，成对方法学习两个文档的先后关系，列表方法则直接面向整个结果列表优化\parencite{liu2009ltr}。学习排序增强了多种信号之间的组合能力，但其效果依赖特征工程和训练数据；点击日志还会受到展示位置、界面设计和用户群体差异的影响，不能被简单视为无偏的相关性标签。

检索效果需要通过与任务目标相匹配的指标评价。准确率（Precision）表示返回结果中相关文档的比例，召回率（Recall）表示全部相关文档中被系统找回的比例，二者分别对应结果纯度与覆盖程度。平均倒数排名（Mean Reciprocal Rank, MRR）重视第一个相关结果出现的位置，适合每个查询通常只需一个正确答案的场景；平均准确率均值（Mean Average Precision, MAP）综合考虑多个相关文档在排序列表中的分布；归一化折损累计增益（Normalized Discounted Cumulative Gain, NDCG）允许设置多级相关性，并对排名靠后的结果给予较小权重\parencite{jarvelin2002cg,manning2008ir}。这些指标从不同侧面衡量排序质量，不能彼此简单替代，也无法完全反映满意度、公平性和实时性等线上体验。

总体而言，候选召回与相关性排序使搜索系统从单纯判断“是否包含查询词”，发展到综合估计“结果是否值得优先展示”。但 BM25 仍以词项统计为基础，链接分析和传统学习排序也往往依赖人工设计的离散特征；面对词汇不匹配、复杂语境和长尾查询时，模型对深层语义关系的表达仍然有限。预训练语言模型随后为查询和文档提供了连续语义表示，推动检索系统进入稠密检索与神经排序阶段。

\subsection{稠密检索与神经排序}

BM25 等稀疏检索方法把查询和文档表示为由词项构成的高维稀疏向量，其优势是索引成熟、匹配过程透明，并能准确处理人名、编号等关键字符串；其不足则在于相关性主要来自相同词项的重合。深度学习与预训练语言模型提供了另一种表示方式：模型将文本编码为低维连续向量，使表述不同但语义接近的查询和文档在向量空间中彼此靠近。检索由此不再完全依赖显式关键词，相关性模型也能够利用上下文区分一词多义、同义改写和自然语言问题。

双编码器是稠密检索中最常见的架构。查询编码器 $E_Q$ 与文档编码器 $E_D$ 分别计算固定维度向量，并以点积或余弦相似度作为得分：
\begin{equation}
  s(q,d)=E_Q(q)^{\mathsf T}E_D(d).
  \label{eq:dense-retrieval}
\end{equation}
由于文档向量可以预先计算，在线查询只需编码一次，并在向量索引中搜索近邻。DPR 使用这一结构在开放域问答数据上训练查询和段落表示，表明仅依靠稠密向量也可以完成大规模候选召回\parencite{karpukhin2020dpr}。实际系统通常采用近似最近邻（Approximate Nearest Neighbor, ANN）算法，避免逐一计算查询与所有文档的相似度。近似搜索提高了速度，却会以一定的召回损失、额外索引内存和构建成本为代价；向量模型更新后，文档编码与索引还可能需要重新生成。

稠密检索的有效性不仅取决于编码器，也取决于训练样本。模型通常利用相关的查询—文档对作为正样本，并通过对比学习拉开负样本的距离。随机负样本往往过于容易，不能提供足够的区分信号；从 BM25 或全库 ANN 结果中选择表面相似但实际不相关的“困难负样本”，能够迫使模型学习更细致的相关性边界。ANCE 通过全库近似近邻检索动态构造困难负样本，说明负样本选择是稠密检索训练的关键环节\parencite{xiong2021ance}。然而，标注错误会把实际相关文档当作负样本，训练数据所属领域和查询分布也会影响模型迁移到新语料时的表现。

双编码器的效率来自查询与文档相互独立的压缩，但单个向量可能丢失局部匹配细节。交叉编码器则把查询和候选文档拼接后共同输入 Transformer，使每一层都能建立查询词与文档词之间的注意力交互，再输出相关性得分。基于 BERT 的段落重排序证明了深层交互对相关性判断的价值\parencite{nogueira2019bert}。由于每个查询—文档对都必须单独执行模型推理，交叉编码器无法经济地遍历整个文档库，通常只用于重排 BM25 或稠密检索返回的前若干个候选结果。因此，现代搜索系统往往形成“快速召回—神经重排序”的级联结构，并通过候选数量控制效果与延迟之间的平衡。

双编码器与交叉编码器之间还存在晚交互方案。以 ColBERT 为例，查询和文档仍被分别编码，但系统保留词元级向量，在检索阶段使用轻量的最大相似度运算聚合局部匹配信号\parencite{khattab2020colbert}。这种设计比单向量表示保存了更多细粒度信息，又允许离线预计算文档表示；代价是索引体积、内存访问和在线计算量通常高于单向量稠密检索。由此可见，稀疏表示、单向量稠密表示、晚交互和交叉编码并非简单的替代关系，而是在匹配精度、吞吐量、存储成本和系统复杂度之间选择不同工作点。

稠密检索与神经排序显著缓解了词汇不匹配问题，但也带来新的限制。向量维度缺乏直接的词义解释，使错误结果更难诊断；模型可能忽略罕见实体、数字或否定词等对答案至关重要的精确线索；训练和推理还需要更多算力，并可能继承数据中的偏差。因而实际应用常将 BM25 的精确词项匹配与稠密表示的语义泛化结合，再使用神经模型重排。更重要的是，这一阶段仍以返回相关文档为终点。当大语言模型需要基于检索内容直接生成回答时，系统还必须考虑文档切分、证据组织、上下文冲突和生成事实性，这推动了检索增强生成框架的发展。

\subsection{检索增强生成}

稠密检索和神经排序改善了相关文档的发现能力，但传统搜索仍将阅读、比较和归纳结果的工作留给用户。大语言模型能够直接生成连贯回答，却主要依靠训练时写入参数的知识，难以及时反映外部信息变化，也不容易说明答案依据。检索增强生成（Retrieval-Augmented Generation, RAG）将两者结合：模型一方面保留参数化知识所提供的语言理解与生成能力，另一方面在推理时从可更新的外部语料中检索非参数化知识。Lewis 等人的工作将神经检索器与序列生成模型联合起来，证明外部文档能够支持知识密集型生成任务\parencite{lewis2020rag}。由此，检索结果不再只是供用户点击的终点，而成为生成模型组织回答的证据输入。

典型 RAG 系统首先建立外部知识库。系统从网页、数据库或内部文档中采集内容，完成去重、格式统一、权限过滤和元数据保留，再将长文档切分为可索引的片段。片段过长会混入与问题无关的信息并占用上下文窗口，片段过短则可能破坏完整语义，使回答所需的条件与结论分散在不同位置；因此通常需要结合标题、段落边界和重叠窗口确定粒度。处理后的片段可写入 BM25 倒排索引、稠密向量索引或两者组成的混合索引，同时保留来源、时间和文档位置，以便生成阶段引用和核验。知识库更新也必须同步索引，否则 RAG 虽然具备外部知识接口，实际获得的信息仍可能过时。

在线阶段从查询处理开始。系统可以对用户问题进行消歧、改写或多步分解，再分别执行稀疏或稠密召回；初始候选经过交叉编码器重排序、重复内容合并和可信来源过滤后，形成供生成器使用的证据集合。最终上下文需要在长度限制内安排片段顺序，并明确区分用户指令、系统要求与外部材料。生成器据此综合多个片段，形成回答并尽可能给出出处。Fusion-in-Decoder 将每个检索片段分别编码，再在解码阶段融合其信息，展示了生成模型聚合多篇证据的能力\parencite{izacard2021fid}。这一过程也说明，RAG 不是在提示词前简单拼接若干搜索结果，而是一条包含数据、检索、排序、上下文构造和生成的完整流水线，其中任一环节的错误都会传播到最终答案。

RAG 的首要风险来自检索质量。召回不足会使关键证据根本无法进入上下文，错误或低质量片段则可能诱导模型给出貌似合理但事实错误的回答。盲目增大检索深度 $k$ 并不能稳定解决这一问题：更多片段会增加证据覆盖，也会引入重复、噪声和相互矛盾的陈述。即使上下文窗口足够长，模型对不同位置的信息利用也并不均衡；相关证据处于长上下文中部时，模型表现可能明显下降\parencite{liu2024lost}。因此，检索深度、重排序阈值、片段顺序和上下文压缩需要联合设计，而不能只追求召回数量。

知识冲突进一步暴露了检索与生成之间的关系。当多个来源说法不一致，或外部证据与模型参数记忆冲突时，生成器可能无依据地偏向其中一方、混合不相容信息，甚至忽略正确证据。RAG 因而只能降低幻觉风险，不能自动保证事实性。可靠系统应比较来源时间与权威性，显式呈现无法消解的分歧，并在证据不足时拒绝作答，而不是要求模型始终生成确定答案。固定地为每个请求检索同样数量的文档也会造成不必要的延迟；Self-RAG 等研究尝试让模型判断何时需要检索，并对证据相关性和自身输出进行反思\parencite{asai2024selfrag}，体现了 RAG 从固定流水线向自适应检索的扩展。

RAG 的评价必须同时覆盖检索与生成。检索阶段可继续使用 Recall@$k$、MRR 或 NDCG 衡量相关证据能否进入靠前位置；生成阶段则需要区分答案是否回应问题、是否与参考事实一致，以及每项陈述是否能够由给定证据支持。RAGAs 将评价拆分为上下文相关性、回答相关性和忠实性等维度\parencite{es2024ragas}，有助于定位失败来自检索器还是生成器，但自动指标和大模型裁判本身仍可能存在偏差，关键应用仍需要人工核验。系统评价还应记录端到端延迟、向量索引存储、重排序与生成调用成本，因为更深的检索和更长的上下文通常意味着更高资源消耗。

总体而言，RAG 把搜索系统的任务由返回相关文档推进为选择、组织并利用可验证证据，也使传统检索中的召回与排序错误直接影响生成质量。随着知识载体从文本扩展到图像、音频和视频，检索器还需在不同模态之间建立对应关系，生成器则要定位并整合细粒度信息；当模型能够自主改写查询、选择工具并决定是否继续搜索时，RAG 又进一步成为智能体与 Web 交互的基础。这些变化构成 WWW 2026 相关研究的重要背景。

\begin{figure}[htbp]
  \centering
  \begin{tikzpicture}[
    node distance=3mm,
    stage/.style={
      draw=black!55,
      rounded corners=1.5mm,
      fill=black!3,
      align=center,
      minimum height=1.65cm,
      text width=2.45cm,
      font=\small
    },
    flow/.style={-{Latex[length=2mm]}, semithick, draw=black!65}
  ]
    \node[stage] (keyword) {关键词检索\\倒排索引};
    \node[stage, right=of keyword] (ranking) {候选召回\\相关性排序};
    \node[stage, right=of ranking] (dense) {稠密检索\\神经重排序};
    \node[stage, right=of dense] (rag) {检索增强生成\\证据组织};
    \node[stage, right=of rag] (agent) {多模态 RAG\\智能体搜索};
    \draw[flow] (keyword) -- (ranking);
    \draw[flow] (ranking) -- (dense);
    \draw[flow] (dense) -- (rag);
    \draw[flow] (rag) -- (agent);
  \end{tikzpicture}
  \caption{Web 搜索与检索增强生成的技术演进}
  \label{fig:retrieval-evolution}
\end{figure}

\subsection{WWW 2026 的研究重点}

WWW 2026 对搜索与检索增强人工智能的征稿范围同时涵盖 Web 爬取与索引、查询理解、搜索与排序、评价方法、垂直搜索、大语言模型搜索、RAG、多模态 RAG 和智能体搜索，其目标是连接传统信息检索与现代人工智能\parencite{www2026track}。官方录用页面在该方向列出 60 篇论文\parencite{www2026accepted}。根据这些论文的题名和研究对象，本文将其主要关注点归纳为表~\ref{tab:research-directions} 所示的五类。该分类并非会议设置的下级类别，而是为了观察研究问题如何沿前述技术主线继续演进。

\begin{table}[htbp]
  \centering
  \caption{本文归纳的 WWW 2026 搜索与检索增强人工智能论文研究重点}
  \label{tab:research-directions}
  \begin{tabularx}{\textwidth}{
    >{\raggedright\arraybackslash}p{0.18\textwidth}
    >{\raggedright\arraybackslash}X
    >{\raggedright\arraybackslash}X}
    \toprule
    研究重点 & 主要问题 & 代表论文 \\
    \midrule
    检索与排序 & 复杂查询下的召回、相关性和扩展效率 & \emph{Generalized Pseudo-Relevance Feedback}；\mbox{\emph{LongRanker}} \\
    高效可靠 RAG & 检索深度、噪声、冲突、幻觉与计算成本 & \emph{PruneRAG}；\emph{CompactRAG}；\emph{Conflict-Aware RAG} \\
    多模态 RAG & 文本、图像、音频和视频的对齐与细粒度证据 & \emph{LoVR}；\emph{IRAG}；\emph{CFVBench} \\
    智能体搜索 & 搜索时机、工具选择、停止条件和轨迹评价 & \emph{DeepAgent}；\emph{TRACE}；\emph{To Search or Not to Search} \\
    垂直应用 & 医学、学术检索、购物与知识图谱 & \emph{Med-R$^2$}；\emph{PaperAsk}；\emph{AgenticShop} \\
    \bottomrule
  \end{tabularx}
\end{table}

首先，传统检索问题并未因生成式人工智能出现而失去意义，而是在新的模型条件下继续发展。录用工作既推广了相关性反馈方法，也探索利用长上下文大语言模型进行一次性文档重排序，并涉及近似最近邻搜索、生成式检索以及科学文档的语义匹配。这些研究覆盖索引、召回与排序的多个环节，说明 RAG 的上限仍受基础检索质量制约。大语言模型在这里既是被搜索系统服务的生成器，也成为查询理解和排序建模的新工具；研究问题并未从“检索”跳跃到“生成”，而是形成了两者相互依赖的系统。

其次，RAG 研究的重点明显从“能否检索到相关内容”扩展为“证据是否可靠且能否被正确使用”。\emph{Rethinking the Hidden Risk of Reranking} 直接关注重排序可能给 RAG 带来的风险，\emph{NeocorRAG} 强调减少无关信息并构建显式证据链，\emph{Conflict-Aware RAG} 处理相互冲突的检索内容，\emph{SeaRAG} 则以自适应排序降低幻觉。与此同时，\emph{PaperAsk}、学术引用预测基准和事实核验研究把可靠性细分为证据相关、引用适当、答案忠实与结论可核验等维度。由此可见，面向生成的检索不能只优化 Recall@$k$ 或 NDCG，还必须判断片段能否支撑具体陈述，并分析错误会如何从重排序和上下文组织传播到最终回答。

第三，效率开始被作为与准确性并列的目标。\emph{PruneRAG} 通过带置信度的查询分解树控制检索过程，\emph{Less is More} 研究紧凑线索选择，\emph{CompactRAG} 减少多跳问答中的模型调用和词元开销，另有工作直接研究具有自适应检索深度的成本感知推理，以及面向查询的上下文压缩。录用列表还包含对大语言模型 Web 搜索引擎实时抓取可持续性的测量研究。这些工作共同指出，增加检索轮数、候选数量和上下文长度虽然可能带来更多证据，却不必然提高答案质量，并会线性乃至更快地推高搜索、重排序与生成成本。因此，当前问题已由固定的“检索前 $k$ 项”转向根据问题难度、证据充分性和资源预算动态决定检索深度。

第四，多模态检索使“文档”概念扩展为带有空间、时间和结构的信息单元。\emph{LoVR} 面向多模态环境中的长视频检索，\emph{IRAG} 关注多模态 RAG 的鲁棒性，相关工作还覆盖组合图像检索、文本行人检索、草图检索、多模态文档解析和带噪视频检索。文本页面通常可以依靠段落边界切分，而视频问题可能需要同时利用语音、画面、字幕和时间位置；相关片段即使被召回，生成模型也未必能够捕捉只出现数帧的数字、图标或界面变化。因此，多模态 RAG 的瓶颈不再只是全局语义是否相似，还包括不同模态能否对齐、证据能否在时间轴上精确定位，以及生成器能否保留细粒度信息。CFVBench 正是在这一问题上提供了集中评价案例。

第五，智能体搜索将检索从一次性模块变为连续决策过程。\emph{To Search or Not to Search} 研究深度搜索智能体的搜索决策边界，\emph{DeepAgent} 和 \emph{ToolBox-RL} 分别关注可扩展工具集下的推理和跨工具库泛化，\emph{TRACE} 则把评价对象扩展到深度研究智能体的完整轨迹。此类系统不仅要回答“检索什么”，还需决定是否搜索、如何改写查询、调用哪个工具、何时停止以及怎样根据中间证据调整计划。医学问答、学术阅读、知识图谱、时间序列预测和个性化购物等应用进一步表明，不同领域对来源权威性、结构化推理、实时性和用户约束具有不同要求，通用检索流程需要与领域知识和任务规则结合。

综合来看，这 60 篇论文延续了关键词检索、索引、召回和排序的基础，同时呈现出四项相互关联的变化：评价目标从相关性扩展到可靠性，系统设计从固定流程转向成本感知的自适应检索，证据类型从文本扩展到多模态内容，搜索行为从单次查询发展为智能体的多步决策。其共同核心是如何为生成模型和智能体提供数量适当、来源可信、位置明确且足以支撑结论的证据。

\subsection{本章小结}

Web 信息获取经历了从关键词匹配与倒排索引，到候选召回和相关性排序，再到稠密检索、神经重排序与 RAG 的演进。技术目标也由高效定位含有查询词的网页，逐步扩展为理解语义、组织外部证据并支持生成可靠回答。WWW 2026 的相关论文进一步显示，可靠性、效率、多模态和智能体化已成为检索研究的重要方向。然而，召回相关内容并不等于生成模型已经正确利用证据：对于长视频，短暂画面、字幕、语音和跨时间片段之间的关系尤其容易丢失。下一章据此分析 CFVBench，考察细粒度视频多模态 RAG 的评价方法、实验发现及其局限。
